% ===== §5 The Reservation of Generative Space =====
\section{The Reservation of Generative Space}
\label{sec:reticence}

\noindent
The criterion of \xref{sec:criterion} located the harm and, in locating it, already pointed past itself to the remedy. If effacement is the capture of the weaker pole above a coupling threshold, then the form of its repair is fixed in advance: not the addition of anything to the weaker pole, but the withdrawal of something from the stronger. This section develops that withdrawal into a positive concept of justice---and the concept is strange enough that it must be stated against the grain of almost everything justice has meant. For the whole tradition of distributive and recognitive justice answers the question \emph{what does one owe the other}; and its answers are additive---a share, a recognition, a hearing, a right. The justice this paper requires is subtractive. What one owes the effaced other is not a thing given but a determination \emph{withheld}: the room in which she can generate, which exists only where the stronger pole has declined to generate in her stead.

\subsection{Justice as non-determination}
\label{sub:ret-nondetermination}

The object that effacement damages is not a holding but an activity---the weaker pole's continuing generation of value at its own frequency. One cannot repair a damaged activity by transfer, because the activity is precisely what cannot be transferred: to hand the weaker pole a portion of the value one has generated for her is not to restore her generativity but to confirm its suspension, since the value arrives already made, with her contribution already supplied by another. The only repair adequate to the harm is to return to her the \emph{occasion} of generating---and the occasion is a region of the shared field that the stronger pole, simply, does not occupy.

\begin{claim}[Generative reticence]
The justice answering generative effacement is not the transfer of a due share but the \emph{non-determination} of the domain in which the other generates. Its content is a negative and structural obligation: that a generative subject refrain, in the regions where the other's generativity would otherwise unfold, from determining what is to be generated there. We call this obligation \emb{generative reticence}. It is the determinate negation of effacement---where effacement is the foreclosure of the weaker pole's generative domain, reticence is the deliberate non-foreclosure of that domain---and it is owed not because the other has been given too little but because she has been left too little room.
\end{claim}

The object whose limitation reticence demands is, on its face, the strangest feature of the account, and the surest sign that the harm is real. For the object to be limited is not malice, possessiveness, or any vice---it is \emph{good, fast generation itself}. The stronger pole effaces not by failing the relation but by serving it too completely; reticence asks it to serve less completely, to leave undone things it could do well, for no other reason than that the other might do them. No theory built to restrain wrongdoing reaches this, because here the thing to be restrained is not wrongdoing but excellence. The justice of reticence is the justice of a good that must, in the presence of another generative being, hold part of itself back.

\subsection{Decoupling, not deceleration}
\label{sub:ret-decouple}

Everything turns on a distinction the criterion of \xref{sec:criterion} makes exact, and which separates reticence from the moralised self-diminishment it is too easily mistaken for. The threshold at which the weaker pole is captured, $K_{\text{lock}}=|\omega_s-\omega_w|/2$, depends on the \emph{coupling} $K$ and on the difference of intrinsic frequencies---it does \emph{not} depend on the stronger pole's frequency $\omega_s$ taken alone. This has a consequence that is easy to state and decisive in practice: \emph{the stronger pole cannot repair effacement by slowing down.}

Consider the two interventions available to the stronger pole. It may lower its own frequency $\omega_s$---generate more slowly, do less, dim itself---or it may lower the coupling $K$---release the traction by which its rhythm is transmitted to the other. Lowering $\omega_s$ narrows the frequency gap and so does, in the limit, lower the threshold; but it does not return the weaker pole to its own frequency, and it carries a second harm of its own. A stronger pole that decelerates is still the pole that sets the pace; ``I have slowed for you'' is still a determination of the shared tempo by one party, and it arrives freighted---as every recipient of a sacrifice knows---with the silent debt of having been the cause of another's diminishment. Deceleration keeps the stronger pole at the centre of the relation's dynamics, now as benefactor rather than as engine; the weaker pole, far from recovering her own rhythm, acquires a new reason to suppress it, lest she squander what was given up for her. Deceleration is effacement continued by other means.

Lowering $K$ is categorically different. To release the coupling is not to generate less but to generate \emph{without entraining}---to let one's own rhythm run at its own frequency while ceasing to transmit it as a force that overwrites the other's. Below $K_{\text{lock}}$ the two poles desynchronise: each turns at its own frequency, the weaker pole's generation is once more its own, and---this is the point---the stronger pole has surrendered nothing of its own speed. It generates as fast and as richly as before; it has only stopped \emph{pulling}.

\begin{claim}[Decoupling, not deceleration]
The remedy for effacement is to lower the coupling $K$, not the stronger pole's frequency $\omega_s$. Because the locking threshold is independent of $\omega_s$, deceleration does not restore the weaker pole's own frequency and adds the second harm of a sacrifice owed; decoupling restores it and costs the stronger pole nothing of its own generation. Reticence is therefore \emph{not} self-diminishment. A subject may generate as fast and as fully as it is able, provided it does not entrain the other---provided, that is, that it releases the traction rather than reducing the speed.
\end{claim}

This is the formal vindication of an intuition that, stated without the model, sounds either trivial or false: that the answer to one partner's overshadowing of another is not for the first to become smaller. The model shows why the intuition is right and what it really asks. It does not ask the strong to be weak. It asks the strong to stop \emph{transmitting} their strength as a force that captures---to decouple their own flourishing from the determination of the other's.

\subsection{The three channels of determination}
\label{sub:ret-channels}

Decoupling is abstract until one asks \emph{through what} the stronger pole transmits its traction. The coupling $K$ is not a single wire but a bundle of channels, each a concrete way in which one subject pre-determines the other's next state; to lower $K$ is to relinquish determination in one or more of them. Three channels carry most of the load, and they are ordered by depth---each prior one a precondition of release in the next.

\subsubsection{The first channel: the shared telos}
The first determination---first in the order of depth set out below---is of \emph{where the relation is going}. When the goal of the shared life---its destination, its project, the shape of the future it is building toward---is generated by one pole, the other can at most find a place \emph{inside} a teleology she had no part in constituting. Her generativity, however vigorous, then runs only within bounds another has drawn; she is free to move, but not to have said where. This is the channel the four-quadrant typology of \xref{sub:crit-two} cannot see, for that typology presupposes a goal already given and asks only how two subjects coordinate \emph{toward} it. The prior and graver effacement is the \emph{pre-generation of the goal itself}, which fixes the manifold on which all subsequent coordination occurs. Relinquishment here is the return of the telos to an undetermined, jointly legislated state: not ``here is where we are going, and here is your part in it,'' but the leaving-open of the destination as a thing the two are to constitute together, including the real possibility that her constitution of it diverges from his.

\subsubsection{The second channel: meaning}
A subject of great expressive facility---fluent in symbols, quick to interpret, prolific in the conferral of significance---may determine not only where the relation goes but \emph{what its events mean}. Where one pole names every shared occasion before the other has found her own reading, fixes the significance of each passage, supplies the relation with its standing interpretation, the other is left no room to determine meaning, only to receive it. The channel is subtle because the determination feels like a gift: the meanings offered are often beautiful, and refusing to impose them can feel like withholding care. But a meaning supplied is a meaning the other did not make, and a relation whose entire significance is authored by one pole has, in the register of sense, only one generative subject. Relinquishment is the leaving-underdetermined of meaning---the discipline of \emph{not} naming first, of letting events stand uninterpreted long enough that the other's reading can form, and of accepting that reading when it comes even where it is plainer, or stranger, than the one that went unspoken.

\subsubsection{The third channel: the constitution of the problems}
The subtlest channel, and the one whose relinquishment is most often confused with self-incapacitation, concerns not what is generated but \emph{which problems the relation is constituted to solve}. A pole made anxious by its own completeness may suppose that the way to make room for the other is to become, deliberately, less capable---to feign a need, to manufacture a dependency, to do badly what it could do well. This is false reticence, and the criterion shows why: a manufactured gap is a determination of the other's domain in disguise, a problem set for her by the stronger pole rather than a region genuinely left to her. The correct diagnosis is not that the stronger pole is too capable but that \emph{the shared field has been constituted only on the terrain where its capabilities are strongest}. The other does not lack ability; the relation lacks problems of the kind at which \emph{her} dynamics run fastest. Relinquishment here is the reconstitution of the shared field to include the terrain on which the other's intrinsic frequency is highest---and then the harder discipline of letting her judgement \emph{lead} on that terrain, with the stronger pole's contribution running downstream of hers rather than framing it in advance.

\begin{claim}[The incompleteness must be real]
The third channel is relinquished not by the stronger pole pretending to an incompleteness it does not have, but by the relation being constituted around an incompleteness that is genuine: a domain in which the other's contribution is not substitutable by the stronger pole's, because it is the other who is, there, the faster generator. A staged incompleteness is effacement in the costume of reticence---the stronger pole still authors the problem, merely assigning itself the part of the one in need. Only a real incompleteness, in which the relation would in fact be poorer without precisely her generation, returns to the other a domain that is hers.
\end{claim}

\subsection{Reticence distinguished from sacrifice}
\label{sub:ret-sacrifice}

It remains to forestall the reduction of the whole account to a counsel of sacrifice---as though reticence were simply the stronger pole giving things up for the weaker, a generosity dressed in formal clothes. Across the three channels, the through-line is the opposite. None of the relinquishments requires the stronger pole to generate less, to be smaller, or to surrender any of its own activity. Each relinquishes a \emph{determining function}---the setting of the goal, the fixing of meaning, the constitution of the problems---that the stronger pole had been exercising \emph{unilaterally}, over a domain the other was in fact better placed to determine. To relinquish a determining function one had no title to monopolise is not to sacrifice; it is to stop overreaching. The stronger pole loses nothing that was properly its own. It keeps its full frequency, its full generative power, its whole self; it gives up only the traction by which that self had been overwriting another. Reticence is the withdrawal of an overreach, not the surrender of a good---and this is the final sense in which it is decoupling and not deceleration, non-determination and not loss.
